% =============================================================================
% Chapter 10 -- Remaining FYP Plan              (owner: Abdullah Amin)
% Derived from TODO.md \S{}3 (deliverables), \S{}4 (dataset), \S{}7 (mid-eval readiness)
% and \S{}9 (frozen --- do not build). Dates are stated relative to the start of
% FYP-2 because no calendar dates have been agreed with the supervisor yet.
% =============================================================================
\chapter{Remaining FYP Plan}
\label{ch:fyp2plan}

This chapter is the plan for what is left. It is written from the team's
working checklist rather than from the proposal, so that the plan, the tasks the
team actually works from, and the report all say the same thing. Every work
package below has one accountable owner; others may help, but one name defends
it at an evaluation.

% ---------------------------------------------------------------------------
\section{Remaining Scope}
\label{sec:remaining-scope}

Of the committed scope in Chapter~1, the reviewer surface, the requirements and
design baseline, the diagram set and the evaluation-data pilot exist. Everything
that turns the system from a described design into a running system remains:
ingestion and reliability, the evidence graph and its rules, and the
verification layer with its gate. Those are three work packages, they are
ordered by dependency, and nothing new has been added to them.

\subsection*{Frozen --- recorded so it cannot creep back}

The following are explicitly \emph{not} planned, from the frozen list in the
team checklist:

\begin{itemize}
  \item New integrations (GitLab, Linear, Notion, Vercel). The adapter layer is
        defined by source role, so each of these would be one adapter against an
        existing contract, and none adds an evaluable claim.
  \item Autonomous code generation or repair, autonomy level L4, two-way
        write-back to source systems, a hosted production LLM, or unrestricted
        chat inside the product.
  \item Pixel-perfect Figma certification and any screenshot, DOM or
        vision-language comparison of design against implementation.
  \item A dedicated graph database, full event sourcing, or Neo4j --- deferred
        until multi-hop pattern queries are actually required.
  \item AI chain summarisation, natural-language graph querying and standup
        summaries: presentation features with no evaluable claim.
  \item Multi-hop cascade analytics and community or centrality analytics.
  \item Learned or adaptive confirmation thresholds, deferred to FYP-II and
        contingent on the fixed-threshold version being stable at the second
        mid evaluation.
\end{itemize}

% ---------------------------------------------------------------------------
\section{Remaining Modules and Features}
\label{sec:remaining-modules}

\begin{longtable}{p{0.16\textwidth} p{0.22\textwidth} p{0.44\textwidth} p{0.14\textwidth}}
\caption{Modules and features remaining for FYP-2, with the planned approach.}\label{tab:remaining-modules}\\
\hline\textbf{Package} & \textbf{Module} & \textbf{Contents} & \textbf{Owner} \\
\hline\endfirsthead
\hline\textbf{Package} & \textbf{Module} & \textbf{Contents} & \textbf{Owner} \\
\hline\endhead
Slice A --- ingestion reliability & \texttt{ingestion/}, \texttt{store/} &
  Mock-vendor replay harness (duplicate, out-of-order, delayed, missing
  webhook, outage, replay); normalisation of the three sources into one
  envelope; idempotency by stable event id; per-resource checkpoints with
  reconciliation polling that synthesises gap events; \texttt{SOURCE\_STALE}
  marking with dependent findings downgraded; immutable artifact and version
  store; the \S{}8.4 reliability state machine.
  & Haaniah \\
Slice B --- code/test evidence graph & \texttt{graph/}, \texttt{analyser/},
  \texttt{conflicts/}, \texttt{resolution/} &
  Postgres property graph (\texttt{nodes}, \texttt{edges},
  \texttt{artifact\_snapshots} with the planned indexes); deterministic linking
  from explicit references plus manual confirm/reject with resolver provenance;
  the link-eligibility gate; TypeScript/React analyser producing base/head diffs,
  changed symbols, a bounded import/call neighbourhood and test evidence;
  symbol cards; requirement anchors with re-anchoring; Rules 1--3 and policy P4
  as versioned records; indexed rule dispatch with reverse-index invalidation;
  role-edge notification routing.
  & Mishal \\
Slice C --- verification and reviewer surface & \texttt{verify/}, \texttt{ui/},
  \texttt{notify/} &
  Deterministic context-packet builder with per-slot budgets and a truncation
  manifest; local model-serving container with pinned model and settings;
  bounded three-pass verifier producing the structured schema; the
  verdict/provenance gate with citation resolution and downgrade reasons; the
  four-verdict semantics with freshness and review fields kept separate;
  binding the existing front end to the API; findings lifecycle and audit;
  Slack delivery adapter (delivery only).
  & Abdullah \\
End-to-end scenarios & \texttt{tests/} (all three) &
  T1--T11 from Chapter~3: explicit-reference chain, Rules 1--3 positive and
  negative, duplicate/out-of-order convergence, outage to
  \texttt{SOURCE\_STALE} to recovery, role routing, fabricated-citation
  downgrade (T10, the demo's failure case), incremental-versus-full
  equivalence, and a degraded-mode run with the model server off.
  & All three \\
\hline\end{longtable}

Three research deliverables from the panel feedback also remain and are
deliverables in their own right: the model-selection study with its resource
budget table (no model may be downloaded before the protocol and the hardware
are recorded); the three-loop analysis showing what the product manager does
versus what the system does, with the intervention metrics that make the claim
measurable; and the completed traceability matrix
(requirement $\rightarrow$ design element $\rightarrow$ module $\rightarrow$ test
case $\rightarrow$ demo step), regenerated at each gate.

% ---------------------------------------------------------------------------
\section{Remaining Iterations}
\label{sec:remaining-iterations}

Development continues in iterations that each terminate in something
measurable, closing at a gate where the named evidence has been reviewed by the
two members who did not produce it.

\begin{longtable}{p{0.14\textwidth} p{0.32\textwidth} p{0.30\textwidth} p{0.20\textwidth}}
\caption{Remaining iterations and the scope assigned to each.}\label{tab:remaining-iterations}\\
\hline\textbf{Iteration} & \textbf{Objective} & \textbf{Gate} & \textbf{State} \\
\hline\endfirsthead
\hline\textbf{Iteration} & \textbf{Objective} & \textbf{Gate} & \textbf{State} \\
\hline\endhead
Iteration 1 (this report) & Requirements and design baseline; reviewer surface;
  evaluation-data pilot & Partially met: the baseline, the surface and the
  pilot exist; the ingestion, linking and spike deliverables did not start
  & Closed with a declared variance \\
Iteration 2 & Ingestion that survives real delivery behaviour; the evidence
  graph and the analyser; criterion decomposition; the context-packet builder;
  the frozen evaluation set & Reliability demonstrated rather than asserted,
  and the evaluation set frozen before any component is tuned against it
  & Not started \\
Iteration 3 & The bounded verifier, the verdict gate, deterministic rules and
  notification routing end to end; the front end bound to the real API &
  Scenario matrix T1--T11 passing, degraded mode demonstrated, hard gates 1 and
  2 enforced & Not started \\
Iteration 4 & Evaluation: benchmark against every baseline, the adaptation
  experiment with its pre-committed go/no-go rule, the crossover workflow study
  & Per-class results reported whether favourable or not; provisional targets
    finalised with evidence & Not started \\
\hline\end{longtable}

% ---------------------------------------------------------------------------
\section{Implementation Plan}
\label{sec:implementation-plan}

Sequenced by dependency, with each item's definition of done fixed in advance:
tests green, an evidence row on the board, the diagram updated if structure
changed, a report note drafted, and any GenAI assistance logged at its level.

\begin{enumerate}
  \item \textbf{Bootstrap the foundation} --- scaffold the remaining workspace
        (\texttt{apps/api}, \texttt{apps/worker},
        \texttt{packages/contracts}, \texttt{db/migrations},
        \texttt{fixtures/}, \texttt{tools/mock-vendor/}); write the Zod
        contracts first (event envelope, node and edge types, rule schema,
        context-packet schema, verdict schema) because the database, the tests,
        the UI and this report all read from them; add
        \texttt{docker-compose.yml} and CI with one other member's review
        required.
  \item \textbf{Slice A} --- normalisation, idempotency, checkpoints,
        reconciliation, staleness and the version store, exercised by the
        replay harness rather than by a live vendor.
  \item \textbf{Slice B} --- the graph schema, deterministic linking with
        provenance, the analyser for the committed language, symbol cards,
        requirement anchors, then the rule set with incremental dispatch.
  \item \textbf{Slice C} --- the packet builder and its manifest, the local model
        container, the three-pass verifier, the gate, and the API binding that
        removes the mock layer from the front end.
  \item \textbf{Integrate} --- the scenario set T1--T11 as end-to-end tests, plus
        the degraded-mode pass with the model server switched off.
\end{enumerate}

The riskiest open question is scheduled as early as it can be: the local
inference spike runs as soon as a contract-stable packet exists, because the
hardware ceiling changes what the rest of Slice C is allowed to promise.

% ---------------------------------------------------------------------------
\section{Testing and Validation Plan}
\label{sec:testing-plan}

\paragraph{Five layers, in order.}
\emph{Data} --- agreement statistics, adjudication, the integrity report and a
dataset card; \emph{Component} --- per-stage metrics measured in isolation
(linking precision, retrieval Recall@k, rule precision, packet byte-identity);
\emph{System} --- an acceptance test per functional requirement from
Chapter~4 plus reliability tests for duplicate, out-of-order and outage
delivery and a security scan; \emph{Comparative} --- baselines including exact
linking, lexical and embedding retrieval, a fixed base-model prompt, a
deterministic-rules-only configuration, a full rescan and the manual workflow;
\emph{Human} --- expert verdict review and the crossover task study against the
manual baseline.

\paragraph{Two hard gates.}
A verdict citing no evidence span is rejected automatically and counted as an
abstention, enforced every time. And the same packet run three times at
temperature 0 must agree at 95\% or the configuration is declared unfit to be
measured, before any other number from it is reported.

\paragraph{Ablations committed in advance.}
Comments and docstrings stripped from symbol cards; path and file-name tokens
removed; structural chunking replaced by fixed windows; the verifier removed so
only the threshold remains; and base versus adapted model under identical
retrieval. Each answers a question the panel asked, and each is reported even if
the answer is that the feature does not help.

\paragraph{Dataset phases still to run.}
Complete the annotation round and compute agreement (gate: Cohen's
$\kappa \geq 0.70$, otherwise revise the manual and re-pilot); adjudicate every
disagreement; expand from the 15-case pilot toward a development set of 20--30
cases and a gold oracle of 30--50, with a training pool of 200--500 for the
adaptation experiment; run the audit and freeze phase (integrity report,
supervisor audit of 10--20\% of the gold set plus all rare classes and all
demonstration cases, split manifest frozen); and only then open the test labels,
after the model, prompt, retrieval and thresholds are frozen. Splits stay
by repository, never by example.

% ---------------------------------------------------------------------------
\section{Deployment Plan}
\label{sec:deployment-plan}

The target footprint has not changed: one self-hosted host running Docker Compose
with the web frontend, the API service, the worker, PostgreSQL, Redis and a local
model-serving container, with read-only connections to the three source systems
and an optional Slack delivery adapter. Two properties are acceptance criteria
rather than aspirations: artifact content does not leave the deployment boundary,
and the deployment must be reproducible by a person who did not write it.

The front end is already deployable as a static bundle today, which is what
makes the mid-evaluation demonstration independent of any back-end availability.
The composition file, the migrations and the container definitions are part of
Slice A's bootstrap and do not exist yet.

% ---------------------------------------------------------------------------
\section{Tentative Timeline}
\label{sec:timeline}

Dates below are relative to the start of FYP-2, because no calendar dates have
been agreed with the supervisor; they are stated as a sequence with rough
proportions rather than as commitments.
\textbf{[TODO: verify --- replace with supervisor-agreed dates before
submission.]}

\begin{longtable}{p{0.20\textwidth} p{0.44\textwidth} p{0.32\textwidth}}
\caption{Tentative FYP-2 timeline of phases and deliverables.}\label{tab:timeline}\\
\hline\textbf{Period (relative)} & \textbf{Work} & \textbf{Checkpoint} \\
\hline\endfirsthead
\hline\textbf{Period (relative)} & \textbf{Work} & \textbf{Checkpoint} \\
\hline\endhead
Weeks 1--2 & Bootstrap, contracts, CI, containers; annotation round completed
  and agreement reported & Foundation reviewable; dataset pilot evidence closed \\
Weeks 3--6 & Slice A end to end with the replay harness; Slice B graph and
  linking & Reliability scenarios T7 and T8 demonstrated \\
Weeks 7--9 & Analyser, symbol cards, anchors, rules; local inference spike &
  Hardware ceiling recorded; retrieval measured \\
Weeks 10--12 & Slice C: packets, verifier, gate; API binding replaces the mock
  layer & T1--T6, T9--T11 passing; T10 demonstrated as the failure case \\
Weeks 13--15 & Dataset expansion and freeze; benchmark against baselines;
  ablations & Frozen oracle; per-class results \\
Weeks 16--18 & Adaptation experiment and its go/no-go decision; user study;
  final report and repository freeze & Evaluation tag; submission \\
\hline\end{longtable}

% ---------------------------------------------------------------------------
\section{Expected Final Deliverables}
\label{sec:deliverables}

\begin{enumerate}
  \item \textbf{A working system} for one workflow: one pull request verified
        against its linked criteria, design version, changed code and tests,
        with an evidence case per criterion and a human resolution recorded.
  \item \textbf{An evaluated system}: benchmark results per class against every
        declared baseline, the ablation table, the two hard gates demonstrated,
        the degraded-mode comparison, and the adaptation experiment reported
        either way under its pre-committed rule.
  \item \textbf{A validated dataset package}: repository catalogue, frozen
        taxonomy and manual, development and gold sets with agreement and
        adjudication records, integrity output, split manifest and dataset card.
  \item \textbf{Evidence of process}: the action register closed or explicitly
        carried, twelve stamped diagrams, the traceability matrix, a task and
        commit record with per-member contributions, the GenAI register and
        prompt journal, and the external dependency table.
  \item \textbf{Reproducibility}: a tagged evaluation commit, a deployment a
        non-author can bring up from the documented steps, and this report in
        its final form with every claim traceable to a commit, test, diagram,
        log or dataset file.
\end{enumerate}
